\section{Introduction}
\label{sec:introduction}

The Markov equation
$$
 x^2+y^2+z^2=3xyz
$$
has a tree of positive integer solutions, and the Frobenius--Markov
uniqueness conjecture asks whether a solution is determined, up to permutation, by
its largest entry. One classical way to encode a solution is by a word:
Christoffel words in two letters parametrize the tree, and under the Cohn
representation the largest entry of a triple becomes a matrix invariant of
the corresponding central word. A hypothetical pair of distinct triples
with the same largest entry would then be a pair of distinct central words
with the same invariant, and one can look for structures in the word
setting that forbid such a coincidence. This article isolates one such
structure and determines it completely.

The structure is a reflection. For the two Cohn matrices
$$
 A=\begin{pmatrix}2&1\\1&1\end{pmatrix},\qquad
 B=\begin{pmatrix}5&2\\2&1\end{pmatrix},
$$
a positive matrix word carries three kinds of information at once: a
literal binary word, an integral Gram matrix, and the arithmetic of a root
of $-1$ modulo its first diagonal entry. Each description makes some
operations transparent and hides others. A reflection acts simply on Gram
roots, but its compatibility with positive words turns out to be far more
restrictive, and the question we answer is exactly how restrictive.

The distinction is already visible in a small example. The palindromes
$baab$ and $abba$ have Gram matrices
$$
 \begin{pmatrix}1130&467\\467&193\end{pmatrix},\qquad
 \begin{pmatrix}1130&437\\437&169\end{pmatrix}.
$$
They have the same label and their normalized roots have midpoint $2/5$.
There is no numerical approximation here. The coincidence itself has been noticed
before, as an equality of matrix entries~\cite[p.~4]{LapointeReutenauer2021} and as
a counterexample to a strengthening of the uniqueness conjecture for the values of
lattice paths~\cite[p.~28]{LagisquetEtAl2021}. Yet their completed letter
counts are $(3,3)$, which rules out central words and hence the usual
Christoffel interpretation. In the Thue--Morse coding of Reutenauer and
Vuillon, recalled in \cref{subsec:thue-morse}, the two palindromes become
$abaaaba$ and $aababaa$, and the equality of their labels becomes the equality
of the lengths of two central words; it is the first such coincidence. The issue is to
determine whether this
failure of primitivity is incidental to the example or forced by the
reflection.

We prove that it is forced, and classify every pair responsible for it.
A palindrome $p$ determines a rational reflection $R_p$, which fixes the
terminal column $e=(2,1)^{\mathsf T}$ and is orthogonal for the Gram form
of $p$. The principal word theorem states that, for all finite words $x$
and $y$,
$$
 [R_p\mu(x)e]=[\mu(y)e]
 \quad\Longleftrightarrow\quad
 x\in\{ab\,p,\ ba\,p\}^*,\quad y=\widehat x,
$$
where brackets denote a projective line and the hat exchanges the
blocks $ab\,p$ and $ba\,p$ (\cref{thm:centre-decoder}). This is a
classification for arbitrary finite words, uniformly in the centre, and
the projective equality always lifts to exact equality of columns. When
$p$ is the central word $g_d$ of a standard directive $d$, the blocks are
the images of $ab$ and $ba$ under the standard morphism $\phi_d$, and the
theorem reads $[R_d\nu_d(x)e]=[\nu_d(y)e]$ if and only if
$x\in\{ab,ba\}^*$ and $y=\widehat x$, with $\nu_d=\mu\circ\phi_d$
(\cref{thm:word-decoder}). For this case we give a second proof by
projective cylinders, which also measures how far a reflected column can
move.

For palindromes, the word theorem has a stronger arithmetic consequence.
Equal labels $M$ and the midpoint condition
$$
 \frac{\rho_X+\rho_Y}{2M}=\frac{\rho_p}{M_p}
$$
are equivalent to a full matrix conjugacy by $R_p$, and the decoder then
forces the two palindromes to be $p\,X$ and $p\,\widehat X$, with $X$ a
concatenation of the blocks (\cref{cor:centre-classification}): no common
prefix needs to be assumed, and the prefix $p$ is forced. If $X$ has $2k$
blocks, the completed counts of both palindromes are $2k+1$ times those of
$p$, so they are not primitive when $k>0$. For a standard ancestor this
gives the preimages $\widetilde w w$ and $\widetilde{\widehat w}\,\widehat w$,
with $w$ a concatenation of $ab$ and $ba$, and completed counts
$(2k+1)L_d(1,1)^{\mathsf T}$ (\cref{thm:palindromic-classification}).
Thus the reflection creates an infinite family of exact palindromic
label equalities while excluding every nontrivial central pair at the
midpoint of any palindrome.

We also retain the information that changes when the midpoint moves. The
conjugacy becomes a two-sided matrix equation with a rank-one term, and
its action on a centered reciprocal cusp coordinate is an affine
reflection. For a hypothetical pair of distinct central words with equal
labels, reduced integral coordinates for the midpoint supply a nonzero
integer defect, and the first-letter cylinders of the standard generators
convert that defect into a quantitative restriction on the common
ancestor. The precise statement is \cref{thm:displaced-ancestor-bound}:
when the two parent traces are comparable, the ancestor label $M_0$ is
bounded by a constant multiple of $M^{1/4}$, the constant being explicit
in the data of the actual pair, with no factorization of $M$ required and
no restriction on the parity of $M$. Along a one-sided directive the two
parents grow apart, the constant degrades, and the bound weakens to a
cubic scale. Neither regime excludes all collisions.

The block exchange around a palindrome $p$ (\cref{prop:any-centre}) has a simple
arithmetic: the label of $p$ divides the common label, and the exchange conjugates
the complementary Gaussian factor (\cref{prop:exchange-arithmetic}). Exchanges
around nested centres compose. When the centre of an exchange is itself one of two
exchanged palindromes, the four words built around them form a square of palindromes
with one label, and $r$ nested levels form a cube with $2^r$ vertices
(\cref{thm:nested-cubes}). The pairs of such a square that are not exchanges have their
root midpoints shifted off every centre, by exactly half the width of a short rung
(\cref{cor:nested-square}). In the
Thue--Morse coding the label is the length of an explicit Christoffel word and the
root its number of letters $b$ (\cref{prop:thue-morse}). An exact enumeration of
all palindromes with label below $10^{26}$ finds no coincidence of labels that the
exchanges do not produce, and every class of equal labels it finds is a single exchange
or a square. This is the content of \cref{conj:structure}, whose proof would imply the
uniqueness conjecture; \cref{sec:outlook} explains how far the reflection goes
towards it. It also measures one consequence of the conjecture: two palindromes with the
same label would have directing words of the same length in the Thue--Morse coding. This
length is half the sum of the partial quotients of the label divided by the root, and for
Markov words the fixed-sum theorem of Lagisquet, Pelantov\'a, Tavenas and
Vuillon~\cite{LagisquetEtAl2021} turns the equality of lengths into an equivalent form of the
uniqueness conjecture (\cref{prop:lengths-uniqueness}). On every reflection axis the
equality holds, since twins there are exchanges.

The prescribed midpoint is essential: a general equal-label pair need not
have it (\cref{rem:quadruple} exhibits one), and away from it the inequalities
above remain necessary conditions rather than exclusions. What the article contributes is the
complete word classification at a specified reflection, together with a
quantitative account of the information that survives off its axis. The
Frobenius--Markov uniqueness conjecture itself is not resolved here.

\subsection{Relation to earlier work}
\label{sec:literature}

The standard morphisms and palindromic extensions used here belong to
classical Sturmian combinatorics. Justin's functional identity relates
iterated palindromic closure to composition of standard
morphisms~\cite[Lemma~2.1]{Justin2005}. Its left and right forms, with the
same conventions used below, appear in Perrin and
Reutenauer~\cite[Section~2]{PerrinReutenauer2023}. The incidence formulas
of Berth\'e, de Luca and Reutenauer describe the completed letter counts
of central words and their standard factorizations~\cite[Section~3]{BertheDeLucaReutenauer2008}.
These results supply the word transport and the primitive-count test in
our application. The involution obtained by reversing a directive in
\cite{BertheDeLucaReutenauer2008} changes the central word by Christoffel
duality; our reflection instead acts on a positive matrix orbit while the
directive remains fixed.

The language that appears in the answer is also familiar. If
$\theta(a)=ab$ and $\theta(b)=ba$, then
$(ab\mid ba)^*=\theta(\{a,b\}^*)$ is the image of the Thue--Morse
substitution. Reutenauer and Vuillon use this substitution together with
palindromic and antipalindromic closure to construct Markov
numbers~\cite{ReutenauerVuillon2017}: in their coding a Markov number is two more than the
length of $\operatorname{Pal}(\theta(\operatorname{Pal}(av)))$, and the injectivity of this
length is equivalent to the uniqueness conjecture~\cite[Theorem~1]{ReutenauerVuillon2017}.
\Cref{prop:thue-morse} identifies that length plus two, the length of the Christoffel word
$a\operatorname{Pal}(\theta(F_a(z)))b$, with our Gram label, for every palindrome $z$.
Abram, Lapointe and Reutenauer
extend the construction to whole Markov triples~\cite{AbramLapointeReutenauer2020}.
Our conclusion concerns the role of this language in a different
problem: it exhausts the finite words whose positive terminal vectors
are carried to another such vector by the prescribed rational reflection.
The proof establishes this exhaustion uniformly for the reflection of every
palindrome, before imposing palindromicity on the reflected words.

Palindromic symmetry is well established in the geometry of two-generator
groups. Kassel and Reutenauer prove that a cyclically reduced basis with
both words of odd length has a unique palindromic conjugate
basis~\cite{KasselReutenauer2007}. Gilman and Keen relate palindromic
words to axes perpendicular to the common perpendicular of the generator
axes~\cite{GilmanKeen2009}. These results concern primitive bases or
geometric axes. The decoder here retains the positive word, the finite
terminal vector and the fixed rational reflection, and applies to every
finite word, including nonprimitive ones.

The distinction between a full matrix and one of its entries matters
when connecting the result to Markov arithmetic. Reutenauer's
Christoffel correspondence is a bijection with whole proper Markov
triples~\cite[Theorem~1]{Reutenauer2009}. More recently, Veselov identifies
Markov fractions with Cohn matrix indices and describes continued
fractions through a mirrored Conway topograph~\cite[Sections~3--4]{Veselov2026}.
Jouteur, Paris-Romaskevich and Thomas characterize modular involutions
exchanging rational pairs and study Springborn operations, including a
$q$-deformed midpoint formula for Farey neighbors~\cite[Sections~5--7]{JouteurParisRomaskevichThomas2026}.
Those constructions retain a tree position, an arithmetic regularity
condition or a whole rational fraction. They do not identify which
arbitrary positive words survive reflection in the orbit considered here.

There are also injectivity theorems for restricted word languages.
Lapointe and Reutenauer prove that a Cohn matrix entry is strictly
increasing in military order on the language of one fixed Sturmian
sequence~\cite[Corollary~1]{LapointeReutenauer2021}. Our comparison allows
arbitrary palindromes, whose completions need not lie in one such
language. Their Gram labels are positive integers, and are Markov
numbers under the additional Christoffel hypothesis. Equal labels and a
specified root midpoint recover a full Gram congruence, which brings the
arbitrary-word decoder into play. This yields an exact classification of
the reflected pairs and shows why every nonempty pair fails the primitive
completed-count condition. No claim is made that an arbitrary equal-label
pair has the prescribed midpoint.


\subsection{Organization}

\Cref{sec:standard-grams} fixes the literal standard morphisms,
completed counts and Gram normalization, and reads the label and the root in the
Thue--Morse coding. \Cref{sec:decoder} proves the finite-word decoder by separating
projective cylinders and reading two letters at a time. \Cref{sec:midpoint} derives
the palindromic classification and the central-word obstruction, extends the
construction, the decoder and the classification to every palindromic centre, shows
that the midpoint hypothesis cannot be removed, and proves that nested exchanges
produce cubes of twins. \Cref{sec:displaced} treats moving midpoints, explains the
weighted scale that must be retained, and proves the common-ancestor bound.
\Cref{sec:outlook} states the conjecture on twins, explains the ladders in which
they appear, reads the lengths of the directing words, and isolates the remaining
mathematical question.
